\subsection{Math Theory}

\noindent
$\displaystyle\sum_{k=0}^{r} \binom{m}{k} \binom{n}{r-k} = \binom{m+n}{r}$ \\ 
$\displaystyle\sum_{i=r}^{n} \binom{i}{r} = \binom{n+1}{r+1}$ \\
$\displaystyle\sum_{x=0}^{n} \binom{x+m}{m} = \binom{n+m+1}{m+1}$ \\ 
$\displaystyle\sum_{x=0}^{n} x \binom{x+m}{m} = \displaystyle\sum_{x=0}^{n} x \binom{x+m}{x} = \displaystyle\sum_{x=0}^{n} (m+1) \binom{x+m}{m+1} = (m+1) \binom{n+m+1}{m+2}$ \\ 
$\displaystyle\sum_{x=0}^{n} x^2 \binom{x+m}{m} = \displaystyle\sum_{x=1}^{n} (m+1)x \binom{x+m}{m+1} = (m+1) \displaystyle\sum_{x=0}^{n-1} (x+1) \binom{x+m+1}{m+1} = (m+1)(m+2) \binom{n+m+1}{m+3} + (m+1) \binom{n+m+1}{m+2}$

\vspace{5pt}
\noindent \textbf{Canonical Coin System} \\
$c_i = \dfrac{r^i - 1}{r - 1}$은 $\forall r \ge 2$에 대해 canonical coin system임. ex) \{1, 3, 7, 15\}

\vspace{5pt}
\noindent \textbf{Jacobsthal's function} \\
\begin{flushleft}
\noindent
\begin{tabular}{|c|c|c|c|}
\hline
$n$ & $P_n$ & Primorial ($P_n\#$) & Jacobsthal $j(P_n\#)$ \\ \hline
1 & 2 & 2 & 2 \\ \hline
2 & 3 & 6 & 4 \\ \hline
3 & 5 & 30 & 6 \\ \hline
4 & 7 & 210 & 10 \\ \hline
5 & 11 & 2,310 & 14 \\ \hline
6 & 13 & 30,030 & 22 \\ \hline
7 & 17 & 510,510 & 26 \\ \hline
8 & 19 & 9,699,690 & 34 \\ \hline
9 & 23 & 223,092,870 & 40 \\ \hline
10 & 29 & 6,469,693,230 & 46 \\ \hline
\end{tabular}
\end{flushleft}

\vspace{5pt}
\noindent $\phi(n) \approx \dfrac{6}{\pi^2}n \approx 0.6079n$

\vspace{5pt}
\noindent \textbf{Maximal Prime Gaps} \\
\begin{flushleft}
\noindent
\begin{tabular}{|c|c|r|}
\hline
Upper Bound ($N$) & Max Gap $g_n$ & Starting Prime $p_n$ \\ \hline
$10^3$ & 20 & 887 \\ \hline
$10^6$ & 114 & 492,113 \\ \hline
$10^9$ & 282 & 473,266,931 \\ \hline
$10^{12}$ & 906 & 999,999,999,763 \\ \hline
$10^{15}$ & 1096 & 908,460,772,044,959 \\ \hline
$10^{18}$ & 1,476 & 1,425,172,824,437,699,411 \\ \hline
$2^{64}$ & 1,550 & 18,361,375,334,787,046,697 \\ \hline
% TODO: https://en.wikipedia.org/wiki/Prime_gap#Maximal_gaps
\end{tabular}
\end{flushleft}

\vspace{5pt}
\noindent \textbf{Bertrand's Postulate}
$\forall n > 1, \exists p \text{ s.t. } n < p < 2n$

\vspace{5pt}
\noindent \textbf{moebius} \\
$n$이하 square-free 수의 개수 $\displaystyle\sum_{i=1}^{n} \mu[i] \left\lfloor \dfrac{n}{i^2} \right\rfloor$   

\vspace{5pt}
\noindent \textbf{Frobenius' coin problem} \\
서로소인 두 자연수 a, b에 대해 an + bm꼴로 나타낼 수 없는 최대값은 ab - a - b \\
즉, $ab - a - b + 1 = (a-1)(b-1)$이상의 자연수는 전부 $an+bm$꼴로 표현 가능 \\
gcd(a, b) = g라면 (a/g - 1)(b/g - 1)g 이상의 g의 배수는 전부 표현 가능 \\

\vspace{5pt}
\noindent \textbf{Cycle Count Generating Functions} \\
순열 $\sigma \in S_n$의 cycle 개수를 $c(\sigma)$, 교란순열의 집합을 $\mathcal{D}_n$이라 하면
\[
\sum_{\sigma \in S_n} x^{c(\sigma)}
=
x(x+1)(x+2)\cdots(x+n-1),
\]
\[
\sum_{\sigma \in \mathcal{D}_n} x^{c(\sigma)}
=
\sum_{q=0}^{n}(-1)^q\binom{n}{q}x^q
\prod_{i=0}^{n-q-1}(x+i).
\]
\[
\sum_{\sigma \in S_n}(-1)^{c(\sigma)}
=
\begin{cases}
1 & n=0,\\
-1 & n=1,\\
0 & n\ge 2,
\end{cases}
\qquad
\sum_{\sigma \in \mathcal{D}_n}(-1)^{c(\sigma)}
=
1-n.
\]